\chapter{수익률을 어떻게 잴 것인가}

과제가 Sharpe를 말하는 순간 ``대출 하나의 수익률''이 필요해진다. 그런데 데이터에는 수익률 컬럼이 없다. 있는 것은 약정 금리, 월 상환액, 받은 총액, 마지막 납부일뿐이다. 이 장은 그것들로 수익률을 만드는 법과, 정의를 어떻게 고르느냐에 따라 결론이 뒤집히는 이유를 다룬다.

\section{원리금 균등 상환}
\textbf{질문.} 1만 달러를 연 12\%로 36개월 동안 빌리면 매달 얼마를 갚나?

매달 같은 금액 $A$를 내는데, 그 36번의 납입을 오늘 가치로 할인해 더하면 빌린 돈 $P$가 되어야 한다. 월 이자율 $i=r/12$로 할인하면
\[
P=\sum_{t=1}^{T}\frac{A}{(1+i)^t}=A\cdot\frac{1-(1+i)^{-T}}{i}
\quad\Longrightarrow\quad
A=P\,\frac{i}{1-(1+i)^{-T}}.
\]

\begin{workout}[월 상환액]
\Use{원리금 균등 상환식(이 절), 등비수열의 합}
\Step{1} 월 이자율: $i=0.12/12=0.01$. \Concept{연 금리를 월로}
\Step{2} 할인 인자: $(1.01)^{-36}=0.6989$. 따라서 $1-(1.01)^{-36}=0.3011$.
\Step{3} $A=10000\times 0.01/0.3011=332.14$. \Concept{원리금 균등 상환}
\Ans{매달 \$332.14. 데이터의 \ft{installment}가 바로 이 값이다.}
\Verify{$332.14\times 36=11{,}957$. 이자 총액 약 \$1,957로, 원금이 줄어드는 만큼 12\% \ensuremath{\times} 3년(\$3,600)보다 작다.}
\Trap{연 금리 12\%를 그대로 36으로 곱해 이자를 계산한다. 원금이 매달 줄어드는 것을 잊은 것이다.}
\end{workout}

\section{IRR: 현금흐름이 실제로 낸 수익률}
\textbf{질문.} 대출이 중간에 부도 나면 수익률은 몇인가?

받은 돈의 흐름을 $CF_0=-P$(투자), $CF_1,\dots,CF_H$(받은 돈)로 적자. 이 흐름의 현재가치를 0으로 만드는 월 할인율 $m$이 있으면, 그것이 이 투자가 실제로 낸 월 수익률이다.
\[
\sum_{t=0}^{H}\frac{CF_t}{(1+m)^t}=0,\qquad \mathrm{IRR}=(1+m)^{12}-1.
\]
식이 다항식이라 닫힌 해가 없으므로 수치로 푼다. 코드는 \ft{returns.irr}이며 Brent 방법(\ft{brentq})으로 근을 찾는다.

\begin{keybox}[IRR]
돈을 넣고 돌려받은 흐름 전체를 하나의 ``연 이자율''로 요약한 값이다. 정상 상환이면 약정 금리와 같고, 부도 나면 그보다 낮다.
\Intuit{이 통장에 넣었다 뺀 돈의 흐름을, 매년 몇 \% 이자 통장과 같다고 볼 수 있나?}
\end{keybox}

\begin{workout}[정상 상환 대출의 IRR]
\Use{IRR 정의(이 절), 월 복리}
\Step{1} 현금흐름: $CF_0=-10000$, $CF_1=\dots=CF_{36}=332.14$.
\Step{2} $m=0.01$을 넣으면 \S2.1의 식에 의해 현재가치가 정확히 0이다. 따라서 월 IRR은 1\%. \Concept{IRR 정의}
\Step{3} 연환산: $(1.01)^{12}-1=0.1268$. \Concept{월 복리}
\Ans{연 IRR 12.68\%. 약정 금리 12\%는 단리식 표기(APR)이고, IRR은 복리라 조금 크다.}
\Verify{\ft{checks.py} 2번: 월 0.01000, 연 0.12683.}
\end{workout}

\begin{workout}[10개월 내고 부도 난 대출]
\Use{현금흐름 복원(아래), IRR, 손실률}
\Step{1} 상황(예시): 10번 냈고, 15개월째 회수금 \$500이 들어왔다. 받은 총액 $=10\times332.14+500=3{,}821.43$.
\Step{2} 손실률 $=1-3821.43/10000=0.618$. \Concept{LGD 정의, \S5.6}
\Step{3} 현금흐름 $CF_0=-10000$, $CF_{1..10}=332.14$, $CF_{15}=500$의 IRR을 수치로 풀면 연 $-76.8\%$.
\Step{4} 같은 대출의 단순 연환산(우승 논문식): $(3821.43/10000)^{12/36}-1=-27.4\%$.
\Ans{IRR $-76.8\%$, 단순 연환산 $-27.4\%$. 같은 대출인데 정의에 따라 세 배 가까이 다르다.}
\Verify{\ft{checks.py} 3번. 손실률 0.618은 우연히 2013년 학습 표본의 실제 평균 LGD(0.618)와 같다.}
\Trap{부도 대출의 IRR이 $-100\%$에 가까울 것이라 생각한다. 이미 받은 원리금이 있어 그보다 덜 나쁘다. 반대로 단순 연환산은 손실을 3년에 걸쳐 나눠 덜 나빠 보이게 만든다.}
\end{workout}

\paragraph{현금흐름 복원.} 데이터에는 월별 납입 기록이 없다. 그래서 \ft{returns.schedule}은 이렇게 근사한다: 받은 총액을 월 상환액으로 나눈 몫 $k_f$만큼은 매달 정상 납입했고, 나머지 $\text{rem}=\text{받은 총액}-k_f A$는 마지막 납부 달($\tau$)에 한꺼번에 받았다. 조기상환은 큰 rem이 일찍 들어오는 것으로, 부도 후 회수는 작은 rem으로 표현된다. 이 근사는 연구 노트의 가정 목록에 올라 있다.

\section{포트폴리오 IRR과 수수료}
투자자는 대출 하나가 아니라 여러 개를 산다. 같은 달에 발행된 대출들의 현금흐름을 더해 한 포트폴리오로 보고 IRR을 구한다(\ft{returns.pooled_irr}). LC는 투자자가 받은 돈에서 1\% 서비스 수수료를 뗀다.

\begin{workout}[정상 9건 + 부도 1건]
\Use{pooled IRR(이 절), 수수료}
\Step{1} 각 \$10,000, 연 12\%, 36개월. 9건은 정상 상환, 1건은 앞 예제처럼 부도.
\Step{2} 월별 합계: 1--10개월 $10\times332.14$, 11--36개월 $9\times332.14$, 15개월에 +500. 0개월 $-100{,}000$.
\Step{3} 받은 돈에 0.99를 곱한다. \Concept{투자자 수수료 1\%}
\Step{4} IRR을 풀면 연 6.85\%.
\Ans{부도율 10\%가 연 12.68\%의 수익을 6.85\%로 깎았다.}
\Verify{\ft{checks.py} 4번.}
\end{workout}

\section{네 가지 수익률 정의와 결론의 뒤집힘}
\begin{table}[H]\centering\small
\begin{tabularx}{\linewidth}{@{}p{34mm}LL@{}}\toprule
정의 & 계산 & 문제점\\\midrule
수업식 & 정상이면 약정 금리, 부도면 0 & 부도 손실(원금)을 무시해 과대평가\\
단순 연환산 & (받은/대출)$^{12/T}-1$ & 돈이 일찍 돌아온 것(재투자)을 무시, 조기상환을 손해로 셈\\
재투자 & 받은 돈을 국채 금리로 만기까지 굴림 & 재투자 금리 가정 필요\\
코호트 IRR & 발행월 포트폴리오의 현금흐름 IRR, 수수료 1\% & 현금흐름 시점이 근사\\\bottomrule
\end{tabularx}
\caption{수익률 정의 네 가지(\ft{src/lc/returns.py}). 이 프로젝트의 주 정의는 마지막 줄이다.}
\end{table}

\begin{keybox}[결과 C1--C3]
2013--16년 36개월 대출을 \textbf{전부 승인}하면 코호트 IRR 5.73\%로 같은 기간 3년 국채(0.87\%)를 48개 발행월 모두에서 이긴다(C2). 지난 우승 논문은 ``전부 승인 \ensuremath{<} 국채''라고 결론냈는데, 이는 단순 연환산과 최근 발행분의 종결 표본에서 비롯된 것으로 보인다(C3, 잠정).
\Intuit{같은 데이터라도 자를 대는 방법이 다르면 다른 답이 나온다. 그래서 자부터 정해야 한다.}
\end{keybox}

\section{비교 기준: 무위험수익률은 무엇인가}
Sharpe와 초과수익에는 ``아무 위험 없이 벌 수 있는 수익'' $r_f$가 필요하다.
\begin{itemize}[leftmargin=*]
\item \textbf{만기를 맞춘 국채}(FRED GS3, GS5): 36개월 대출에는 3년 국채, 60개월에는 5년 국채를 발행월 기준으로 붙인다(\ft{returns.attach_treasury}). 투자 기간이 같은 대안이라 가장 정직하다.
\item \textbf{3개월 국채}(TB3MS): 교과서의 $r_f$. 만기가 달라 기간 프리미엄이 섞인다.
\item \textbf{회사채 BAA}: 무위험은 아니지만 ``비슷한 위험의 다른 투자''라는 실무적 비교 기준. BAA로 바꾸면 전부 승인의 우위가 사라지고 선별의 가치가 커진다(C8).
\end{itemize}

\begin{insight}
실무에서 대출 포트폴리오의 성과는 조달 비용(은행이면 예금 금리 + 스프레드) 또는 같은 위험의 벤치마크 대비로 본다. ``무위험''은 출발점일 뿐이라, 이 프로젝트는 국채를 주 기준으로 두고 BAA와 +1--3\%p 프리미엄을 민감도로 함께 보고한다.
\end{insight}

\begin{career}
``Accuracy가 아니라 무엇을 최적화했나?''라는 질문에 이 장이 답이다. 라벨(부도 0/1)이 아니라 \textbf{현금흐름에서 복원한 IRR}을 목표로 삼았고, 정의 선택이 결론을 뒤집는다는 것을 수치로 보였다.
\end{career}
